\section*{AI-use statement}

This paper was drafted by a crew of AI research-assistant agents
(Free-Research-Crew, protocol v0.1) under the direction of the human author, Bryce Wang.
He approved the research question (2026-11-19, 01:17 UTC), the design (2026-11-19,
06:52 UTC) and the results (2026-11-19, 08:01 UTC) at three human gates, each in chat after
reading the Reviewer's verdict and kill argument. The manuscript itself is subject to a
fourth, the submission gate. He takes responsibility for the paper's content.

\begin{itemize}
  \item \textbf{Roles and models.} Every crew role ran on Anthropic Claude through Claude
        Code. The PI role was the main Claude Code session. The Scout, Data RA, Designer,
        Econometrician, Figures RA, Reviewer, Replicator and Writer ran as agents launched
        from that session. The review records name the model as
        \texttt{claude-opus-5-5}. The PI's brief to each agent was archived in
        \texttt{briefs/} from the question stage on, beginning with the Scout's first brief
        (2026-11-18, 16:28 UTC).
  \item \textbf{Review.} The Reviewer ran 3 rounds on the research question, 1 on the
        design, 1 on the code, 2 on the results and 2 on the paper (round~1: revise; round~2: pass with minor
        findings open). No separate
        literature round was held. The code review passed before any analysis run; it came after the design gate, not
        before it as the plan's task said (Section~\ref{sec:deviations}). The
        Reviewer's independence level was \emph{fresh-context}: a new context in every
        round, with the same model family as the other roles. The human author chose this
        over a cross-family reviewer. The final round of each stage passed with minor
        findings still open, all answered before the next gate:
        \begin{itemize}
          \item question round~3 (three findings), answered in the design;
          \item design round~1 (six findings), answered in the plan before the design gate;
          \item code round~1 (three findings), fixed before any analysis run;
          \item results round~2 (two findings), answered in the run notes and claims before
                the results gate;
          \item paper round~2 (four findings), answered in this version: an attribution of
                the H2 narrowing to particular stores that no logged run supports was
                removed, the findings of the two results reviews are now assigned to the
                right round, ``every number'' was narrowed to ``every result number'', and
                this item was completed.
        \end{itemize}
  \item \textbf{Findings accepted over the Reviewer's objection.} None.
  \item \textbf{Evidence.} Every result number in the text and tables is a macro generated
        from the metrics of a logged run, or is written by a logged table script from run
        metrics. The numbers typed by hand are of three kinds:
        \begin{itemize}
          \item design facts fixed in the plan before any analysis run: the minimum wages,
                the survey dates, the targets, the 0.05 tolerance, the 5th and 95th
                percentiles, the store and closure counts checked at data preparation, and
                the Reviewer's pre-run calculation for the row-3 New Jersey standard error;
          \item published or source facts, quoted as the Scout's source notes record them,
                and arithmetic on published entries only;
          \item process facts: dates, review-round counts and the file size of the data
                archive.
        \end{itemize}
        Every citation was checked against its source by the Scout. Of the 12 sources cited, 5
        were read in full text (the original article's tables from page images), 2 in part
        (\citet{neumark1995payroll}: abstract and pp.~1--10, the rest keyword-searched;
        \citet{card1998reanalysis}: abstract, introduction and Section~I.A), 3 from the
        abstract only and 2 by resolving their DOI; the paper cites each for no more than
        was read. The full workspace is public at
        \url{https://github.com/brycewang-stanford/Free-Research-Crew/tree/main/gallery/ss-card-krueger}.
        It contains the plan, the run ledger and notes, the code, the data manifest, every
        review, the archived briefs and the source notes.
  \item \textbf{Replication.} Passed. The Replicator's clean rerun, recorded in
        \texttt{replication.json} (2026-11-19, 12:29 UTC), copied the workspace, deleted every recorded output, fetched the data
        again rather than reusing it, and re-executed all 15 logged runs in a fresh
        environment built only from \texttt{code/requirements.txt}. All 15 reproduced
        their logged metrics within a relative tolerance of $10^{-6}$. The steps are in
        \texttt{REPLICATION.md}.
\end{itemize}
